\section*{第 \ref{ch:g12:early-earth} 章——早期地球及其大气}

\begin{solution}{exo:g12:early-earth:1}
（a）\omterm{def:g12:early-earth:primitive-atmosphere}{脱气}；（b）\omterm{def:g12:early-earth:stromatolite}{叠层石}；（c）\omterm{def:g12:early-earth:bif}{条带状铁建造}；（d）\omterm{def:g12:early-earth:red-beds}{红
层}；（e）\omterm{def:g12:early-earth:goe}{大氧化事件}。
\end{solution}

\begin{solution}{exo:g12:early-earth:2}
显示液态水的锆石（44 亿年）；最早的\omterm{def:g12:early-earth:stromatolite}{叠层石}
出现于皮尔巴拉（3.43）；哈默斯利\omterm{def:g12:early-earth:bif}{条带状铁建造}（超
过 2.5）；\omterm{def:g12:early-earth:goe}{大氧化事件}结束（2.06）；大多数\omterm{def:g12:early-earth:bif}{条带状铁
建造}结束（1.8）。
\end{solution}

\begin{solution}{exo:g12:early-earth:3}
当 $t = 0$ 时，$L = L_\odot/(1 + 0.4) = L_\odot/1.4 \approx 0.714\,L_\odot$：
比今天暗约 \qty{29}{\percent}，亮度低近三分之一。
\end{solution}

\begin{solution}{exo:g12:early-earth:4}
当时空气中几乎没有氧气。黄铁矿会被氧气破坏，
因此其颗粒要能被河流搬运并完整沉积，
前提是地表的空气和水中都缺氧。
\end{solution}

\begin{solution}{exo:g12:early-earth:5}
水蒸气凝结成了海洋，另有二氧化碳，
其在大气中的含量远高于今天，通过
温室效应使地球变暖。
\end{solution}

\begin{solution}{exo:g12:early-earth:6}
30 亿年前：$1 - t/t_\odot = 3/4.567 \approx 0.657$，$L \approx
1/1.263 \approx 0.792\,L_\odot$，$T \approx 255 \times 0.792^{1/4} \approx
\qty{241}{K}$。24.5 亿年前：$1 - t/t_\odot \approx 0.536$，$L
\approx 1/1.215 \approx 0.823\,L_\odot$，$T \approx 255 \times 0.823^{1/4}
\approx \qty{243}{K}$。
\end{solution}

\begin{solution}{exo:g12:early-earth:7}
\qty{1}{kg} 铁约为 $1000/55.85 \approx 17.9$ 摩尔；一个氧分子需要
四个铁离子，因此 $17.9/4 \approx 4.48$ 摩尔的
\ce{O2}，即 $4.48 \times 32.00 \approx \qty{143}{g}$ 的氧气。
\end{solution}

\begin{solution}{exo:g12:early-earth:8}
（1）$2.46 - 2.06 = 0.40$，即 4 亿年。（2）从 45.4 亿年前到 24.6 亿年前，
约占地球历史的 $2.08/4.54 \approx \qty{46}{\percent}$。
（3）约 $3.43 - 2.46 \approx 1$，即 10 亿年：
微生物产生的氧气长期被消耗，
此后才得以留存于大气中。
\end{solution}

\begin{solution}{exo:g12:early-earth:9}
\omterm{def:g12:early-earth:bif}{条带状铁建造}表明：\omterm{def:g10:layered-ocean:thermocline}{深海}缺氧，其中可携带溶解的铁，
其上为阳光充足的浅海，氧气足以使铁沉淀。
\omterm{def:g12:early-earth:red-beds}{红层}形成于陆地，表明空气和地表水中都有氧气，
足以使颗粒中的铁生锈。
\end{solution}

\begin{solution}{exo:g12:early-earth:10}
位于澳大利亚西部鲨鱼湾的哈默林池，其海水盐度
为通常海水的两倍。微生物席以蓝细菌为主，
捕获并黏结沉积物颗粒，促使碳酸盐生长；新的微生物席
又在其上生长，如此层层叠加。
\end{solution}

\begin{solution}{exo:g12:early-earth:11}
（1）$21 \times 10^{-5} = \qty{0.00021}{\percent}$。（2）$21/10 =
\qty{2.1}{\percent}$。
\end{solution}

\begin{solution}{exo:g12:early-earth:12}
若 $L(1 - A')$ 等于 $1 - A$，则\omterm{def:g10:habitable-planet:equilibrium-temperature}{平衡温度}不变：
$1 - A' = 0.706/0.741 \approx 0.953$，故 $A' \approx 0.047$。地球需
只反射不到 \qty{5}{\percent} 的太阳光，几乎如同
黑体，而有云和部分冰层的海洋行星无法做到。
另一种途径是更强的温室效应。
\end{solution}

\begin{solution}{exo:g12:early-earth:13}
（1）升至凝固点需 $273 - 237 = \qty{36}{K}$；升至今天的平均温度需 $288 - 237 = \qty{51}{K}$。
（2）今天的温室效应使温度升高 $288 - 255 =
\qty{33}{K}$，年轻的地球需要更强的温室效应。（3）\qty{0.1}{bar} 相当于
今天整个大气压强的十分之一，且仅来自二氧化碳。
\end{solution}

\begin{solution}{exo:g12:early-earth:14}
（1）早于 24.5 亿年前，该特征存在，说明当时大气中的氧含量
不足今天氧含量的 $10^{-5}$；晚于 20.9 亿年前，该特征消失，说明大气中的氧含量
较多。（2）是的：\omterm{def:g12:early-earth:goe}{大氧化事件}约始于 24.6 亿年前，
止于 20.6 亿年前，二者时段相同。
\end{solution}

\begin{solution}{exo:g12:early-earth:15}
（1）约 $4567 - 60 = 4507$，即 45.07 亿年前。（2）至多 $4507 - 4404 =
103$，即撞击后至多 1.03 亿年：地球在撞击后
约一亿年内就已具有地壳和水。
\end{solution}

\begin{solution}{pb:g12:early-earth:1}
\textbf{1.} $4.567 - 4.0 = 0.567$，即 5.67 亿年。

\textbf{2.} $1 - 0.567/4.567 \approx 1 - 0.124 = 0.876$。

\textbf{3.} $L = 1/(1 + 0.4 \times 0.876) = 1/1.350 \approx 0.741$，
为今天的值。

\textbf{4.} $0.741 \times 1361 \approx \qty{1008}{W/m^2}$。

\textbf{5.} $1 - 0.741 \approx 0.26$：比今天暗 \qty{26}{\percent}。

\textbf{6.} $T_{\mathrm{eq}} = \left(S(1-A)/(4\sigma)\right)^{1/4}$。

\textbf{7.} 当 $A$ 与 $\sigma$ 固定时，$T_{\mathrm{eq}}$ 为常数乘以
$S^{1/4}$，故 $T'/T = (S'/S)^{1/4}$。

\textbf{8.} $0.741^{1/4} = \sqrt{\sqrt{0.741}} \approx \sqrt{0.861} \approx
0.928$。

\textbf{9.} $255 \times 0.928 \approx \qty{237}{K}$。

\textbf{10.} $237 - 273 \approx \qty{-36}{\celsius}$。

\textbf{11.} 地表曾存在液态水，即一片海洋，
距今 44 亿年前。

\textbf{12.} 在浅水中，与今天的活\omterm{def:g12:early-earth:stromatolite}{叠层石}一样。

\textbf{13.} 哈默斯利山脉的\omterm{def:g12:early-earth:bif}{条带状铁建造}，
沉积于海底。

\textbf{14.} 是：\qty{237}{K} 远低于 \qty{273}{K}；若保持今天的
大气，地表平均将处于凝固状态。

\textbf{15.} 岩石表明，液态水存在于一个过于黯淡的太阳之下，
若大气与今天相同，这样的太阳无法维持水的液态：即\omterm{def:g12:early-earth:faint-young-sun}{黯淡太阳悖论}。

\textbf{16.} $288 - 255 = \qty{33}{K}$。

\textbf{17.} 至少 $273 - 237 = \qty{36}{K}$，多于今天的数值。

\textbf{18.} $A' = 1 - 0.706/0.741 \approx 0.047$：几乎全黑的行星，
不合理。

\textbf{19.} 二氧化碳含量远高于今天；甲烷由
早期微生物产生，也可能起了作用。

\textbf{20.} 地球在 40 亿年前的\omterm{def:g10:habitable-planet:equilibrium-temperature}{平衡温度}为
\textbf{约 237 K}，比凝固点低约 \qty{36}{\celsius}。
\end{solution}
